\documentclass[12pt]{beamer}
\usetheme{Madrid}
\usecolortheme{beaver}
\usepackage[utf8]{inputenc}
\usepackage[spanish]{babel}
\usepackage{graphicx}
\usepackage{booktabs}
\usepackage{amsmath}
\usepackage{multirow}
\usepackage{siunitx}
\usepackage{hyperref}

% Configuración de metadatos
\title[Redes Complejas: Cáncer vs Normal]{Análisis Comparativo de Redes Complejas:\\Cáncer de Mama vs. Células Normales}
\author[A. Guerra González]{Alejandro Guerra González}
\institute[FCQ-UASLP]{
	Maestría en Ciencia en Ingeniería Química\\
	Facultad de Ciencias Químicas\\
	Universidad Autónoma de San Luis Potosí
}
\date{\today}

\begin{document}
	
	% Diapositiva de título
	\begin{frame}
		\titlepage
	\end{frame}
	
	% Índice
	\begin{frame}{Índice}
		\tableofcontents
	\end{frame}
	
	\section{Introducción}
	\begin{frame}{Introducción}
		\begin{itemize}
			\item La biología de sistemas transforma la visión reduccionista hacia un enfoque integrador basado en redes de interacción.
			\item Los sistemas biológicos operan como estructuras complejas donde múltiples biomoléculas interactúan coordinadamente para generar fenotipos.
			\item En cáncer de mama, el análisis de redes de interacción génica permite identificar alteraciones estructurales y funcionales entre células normales y tumorales.
			\item \textbf{Objetivo:} Construir y comparar dos redes complejas (normal vs. cáncer) para identificar patrones topológicos y funcionales emergentes asociados a la patología.
		\end{itemize}
	\end{frame}
	
	\section{Métodos y Teoría}
	\begin{frame}{Métodos y Teoría}
		\begin{itemize}
			\item \textbf{Fundamento:} Grafos con nodos (genes/proteínas) y aristas (interacciones). Redes \textit{scale-free} con organización modular.
			\item \textbf{Integración de datos:} Cruce de RegulonDB, BioGRID, KEGG, HPRD, STRING e IntNetDB para aumentar cobertura y reducir ruido experimental.
			\item \textbf{Modelo Bayesiano:} Integración mediante Bayes Naïve. Cálculo de \textit{Likelihood Ratios} (LR) individuales y multiplicación para obtener LR total de confianza.
			\item \textbf{Análisis de redes:} Filtrado por umbral de probabilidad. Métricas: coeficiente de clustering, longitud media de camino, modularidad e identificación de \textit{hubs}.
		\end{itemize}
	\end{frame}
	
	\section{Estructura Topológica Global}
	\begin{frame}{Estructura Topológica Global}
		\begin{columns}
			\begin{column}{0.48\textwidth}
				\textbf{Red Normal}
				\begin{itemize}
					\item Arquitectura dispersa y modular ($Q = 0.7993$).
					\item Módulos compartimentados: UPR, señalización hormonal, metabolismo lipídico.
					\item Fronteras funcionales bien delimitadas.
				\end{itemize}
			\end{column}
			\begin{column}{0.48\textwidth}
				\textbf{Red Cáncer de Mama}
				\begin{itemize}
					\item Componente gigante central denso e hiperconectado.
					\item Hubs oncogénicos: AKT1, ERBB2, PIK3CA, MYC, CTNNB1.
					\item Colapso de programas funcionales en estructura autorreforzante.
					\item Pérdida de identidad tisular.
				\end{itemize}
			\end{column}
		\end{columns}
		\vspace{0.3cm}
		\centering \footnotesize \textit{(Insertar Figuras 1 y 2 del reporte original)}
	\end{frame}
	
	\begin{frame}{Nodos, Aristas y Densidad}
		\begin{itemize}
			\item \textbf{Nodos:} 62 en ambas redes (controla tamaño muestral).
			\item \textbf{Aristas:} 77 (cáncer) vs. 62 (normal) $\Rightarrow$ Incremento del $\sim 24\%$.
			\item \textbf{Densidad ($\rho$):} $0.0407$ (cáncer) vs. $0.0328$ (normal).
			\item \textbf{Interpretación:} Mayor compactación y entrelazamiento de rutas en el contexto tumoral, coherente con hiperactivación de vías de señalización y pérdida de especificidad funcional.
		\end{itemize}
		\begin{equation*}
			\rho = \frac{2m}{n(n-1)}
		\end{equation*}
	\end{frame}
	
	\section{Dinámica de Propagación e Integración de Señales}
	\begin{frame}{Dinámica de Propagación e Integración de Señales}
		\begin{columns}
			\begin{column}{0.5\textwidth}
				\textbf{Coeficiente de Clustering ($\bar{C}$)}
				\begin{itemize}
					\item Cáncer: $0.2716$
					\item Normal: $0.0645$
					\item $\Rightarrow$ Diferencia de $>4\times$.
				\end{itemize}
				\small{Formación de módulos redundantes y bucles de retroalimentación positiva. Alta robustez frente a perturbaciones terapéuticas.}
			\end{column}
			\begin{column}{0.5\textwidth}
				\textbf{Longitud Promedio de Camino ($\bar{\ell}$)}
				\begin{itemize}
					\item Cáncer: $4.75$
					\item Normal: $6.34$
				\end{itemize}
				\small{Red normal: arquitectura jerárquica con filtrado/validación de señales. Red tumoral: propagación rápida y sin restricciones de señales proliferativas.}
			\end{column}
		\end{columns}
		\begin{equation*}
			C_i = \frac{2t_i}{k_i(k_i-1)}
		\end{equation*}
	\end{frame}
	
	\section{Organización Modular}
	\begin{frame}{Organización Modular}
		\begin{itemize}
			\item \textbf{Modularidad ($Q$):} Mide la fuerza de división de una red en comunidades funcionales.
			\item Normal: $Q = 0.7993$ (comunidades claramente delimitadas).
			\item Cáncer: $Q = 0.6689$ (erosión de fronteras funcionales).
			\item \textbf{Significado biológico:} La menor modularidad en cáncer refleja desdiferenciación celular y capacidad del tumor para reclutar y subvertir múltiples rutas simultáneamente.
		\end{itemize}
		\begin{equation*}
			Q = \frac{1}{2m} \sum_{i,j} \left[ A_{ij} - \frac{k_i k_j}{2m} \right] \delta(c_i, c_j)
		\end{equation*}
		\vspace{0.2cm}
		\centering \footnotesize \textit{(Insertar Diagrama de Radar - Figura 4)}
	\end{frame}
	
	\section{Roles Funcionales de los Genes}
	\begin{frame}{Roles Funcionales de los Genes}
		\begin{itemize}
			\item \textbf{Reparación de ADN (11 genes):} Exclusivo del cáncer. Refleja desregulación/mutación en contexto de inestabilidad genómica, no mejora reparadora.
			\item \textbf{Oncogenes (11 genes):} Predominan en cáncer. Actúan como \textit{hubs} que concentran y amplifican señales proliferativas.
			\item \textbf{Other (29 normal vs 8 cáncer):} Mayor diversidad funcional en tejido sano. Contracción en tumor indica pérdida de identidad tisular.
			\item \textbf{ECM/Epitelial y Lactación:} Ausentes o drásticamente reducidos en cáncer. Consistente con transición epitelio-mesenquimal (TEM) y pérdida de diferenciación terminal.
		\end{itemize}
		\centering \footnotesize \textit{(Insertar Tabla 2 y Figura 5 del reporte)}
	\end{frame}
	
	\section{Síntesis Comparativa}
	\begin{frame}{Síntesis Comparativa}
		\begin{table}
			\centering
			\small
			\begin{tabular}{lcc}
				\toprule
				\textbf{Propiedad} & \textbf{Cáncer de Mama} & \textbf{Células Normales} \\
				\midrule
				Aristas & 77 (mayor) & 62 (menor) \\
				Densidad & 0.0407 & 0.0328 \\
				Clustering & 0.2716 & 0.0645 \\
				Long. camino & 4.75 & 6.34 \\
				Modularidad & 0.6689 & 0.7993 \\
				Señal dominante & Proliferación & Especialización \\
				Arquitectura & Redundante & Jerárquica \\
				\bottomrule
			\end{tabular}
			\caption{Comparación de propiedades estructurales de ambas redes.}
		\end{table}
	\end{frame}
	
	\section{Implicaciones Terapéuticas}
	\begin{frame}{Implicaciones Terapéuticas}
		\begin{enumerate}
			\item \textbf{Terapias dirigidas a hubs:} Atacar nodos centrales o desestabilizar cliques es más eficaz que enfoques de blanco único.
			\item \textbf{Resistencia por rutas alternativas:} Baja modularidad y alto clustering justifican terapias combinadas que apunten a múltiples nodos del mismo módulo.
			\item \textbf{Restauración de modularidad:} Estrategias que induzcan diferenciación podrían complementar terapias citotóxicas, abordando la causa estructural de la desregulación.
			\item \textbf{Marcadores topológicos:} Clustering y modularidad como biomarcadores de red para monitorizar progresión tumoral o respuesta terapéutica en estudios longitudinales.
		\end{enumerate}
	\end{frame}
	
	\section{Conclusiones}
	\begin{frame}{Conclusiones}
		\begin{itemize}
			\item La transformación maligna conlleva una reorganización estructural profunda de la red de interacciones génicas.
			\item La red tumoral es más densa ($+24\%$ aristas) e internamente más interconectada ($\times 4$ clustering), favoreciendo propagación incontrolada de señales oncogénicas.
			\item La red normal mantiene una arquitectura de control jerárquica y especializada, perdida en el contexto tumoral.
			\item La distribución de roles confirma dominancia de oncogenes/reparación aberrante en cáncer, y pérdida de identidad epitelial/especializada.
			\item Las diferencias estructurales justifican estrategias terapéuticas basadas en topología de red, orientadas a desestabilizar módulos patológicos en lugar de atacar blancos aislados.
		\end{itemize}
	\end{frame}
	
	\section{Referencias}
	\begin{frame}[allowframebreaks]{Referencias Seleccionadas}
		\footnotesize
		\begin{thebibliography}{10}
			\bibitem{albert2005} Albert, R. Scale-free networks in cell biology. \textit{J. Cell Sci.}, 118(21):4947–4957, 2005.
			\bibitem{barabasi2011} Barabási, A.L., Gulbahce, N., Loscalzo, J. Network medicine. \textit{Nat. Rev. Genet.}, 12:56–68, 2011.
			\bibitem{barabasi2004} Barabási, A.L., Oltvai, Z.N. Network biology. \textit{Nat. Rev. Genet.}, 5:101–113, 2004.
			\bibitem{baranwal2016} Baranwal, M. et al. Network topology measures for identifying disease-gene association. \textit{BMC Bioinformatics}, 17:349, 2016.
			\bibitem{bertho2021} Bertho, M. et al. Activation of PI3K/AKT/mTOR pathway causes drug resistance. \textit{Front. Pharmacol.}, 12:628690, 2021.
			\bibitem{bhattacharyya2014} Bhattacharyya, M., Bandyopadhyay, S. Two-layer modular analysis. \textit{BMC Syst. Biol.}, 8:81, 2014.
			\bibitem{chakrabarti2012} Chakrabarti, R. et al. ELF5 inhibits EMT. \textit{Nat. Cell Biol.}, 14:1212–1222, 2012.
			\bibitem{deng2006} Deng, C.X. BRCA1: cell cycle checkpoint, genetic instability. \textit{Nucleic Acids Res.}, 34(5):1416–1426, 2006.
			\bibitem{dongre2019} Dongre, A., Weinberg, R.A. New insights into EMT. \textit{Nat. Rev. Mol. Cell Biol.}, 20:69–84, 2019.
			\bibitem{fontemaggi2014} Fontemaggi, G. et al. Gain of function mutant p53. \textit{Oncotarget}, 5(4):1006–1022, 2014.
		\end{thebibliography}
	\end{frame}
	
	% Diapositiva final
	\begin{frame}
		\centering
		\Huge ¡Gracias por su atención!
		\vspace{0.5cm}
		\Large Preguntas y Comentarios
	\end{frame}
	
\end{document}